\graphicspath{{Appendici/A01-statica_grafica/Figure/}}
\newcommand{\FiguraSGR}[1]{\input{Appendici/A01-statica_grafica/Figure/#1}}

\chapter{Elementi di Statica Grafica}
\label{app:statica_grafica}

\begin{sommario}
	Si presentano le costruzioni grafiche con cui si riducono e si equilibrano i sistemi piani di
	forze. Il poligono delle forze fornisce la risultante; il poligono funicolare, ottenuto con una
	successione di operazioni elementari di equivalenza, ne individua la retta d'azione, cioè l'asse
	di sollecitazione, e misura il momento del sistema. Letto come catena di aste in equilibrio, il
	poligono funicolare introduce la notazione dei campi, che mette in corrispondenza il piano di
	posizione con il piano delle forze. Se ne ricavano le condizioni grafiche di equilibrio e le
	costruzioni delle reazioni vincolari: il lato di chiusura per le forze parallele, il concorso
	di tre forze per il corpo vincolato da una cerniera e da un carrello e per i sistemi a tre
	cerniere. La stessa notazione conduce all'equilibrio grafico del nodo canonico e al diagramma
	cremoniano delle travature reticolari. Si mostra infine che queste costruzioni hanno un esatto
	corrispettivo nella cinematica grafica dei centri di rotazione.
\end{sommario}

\section{Considerazioni introduttive}
\label{sec:intro_sgr}

La statica grafica tratta per via geometrica i sistemi di forze: le forze sono rappresentate da
segmenti orientati, e le operazioni su di esse da costruzioni eseguite con riga e squadra. Come
procedimento di calcolo è stata in larga misura sostituita dai metodi analitici; conserva invece
intatto il suo valore come rappresentazione e come controllo, perché raccoglie in un solo disegno
la geometria di una struttura e le forze che la percorrono, e con i sistemi di disegno assistito
le sue costruzioni sono tornate a essere esatte e immediate.

L'appendice non aggiunge risultati a quelli del volume: li ritrova per altra via. Presuppone la
statica dei sistemi di forze del \CAPref{cap:statica_CR}, della quale usa la risultante, il
momento risultante, l'asse di sollecitazione (\PARref{sec:asse_soll}) e le operazioni elementari
di equivalenza (\PARref{sec:equivalenza}), e applica le costruzioni a strutture già risolte per
via analitica nei capitoli, così che ogni disegno abbia la propria controprova.

Ogni costruzione si svolge su due disegni distinti, ciascuno con la propria scala: il \emph{piano
di posizione}, nel quale sono tracciate la struttura e le rette d'azione delle forze e si leggono
le lunghezze, e il \emph{piano delle forze}, nel quale le forze sono riportate come vettori liberi
e se ne leggono l'intensità, la direzione e il verso. Dall'uno all'altro si trasportano soltanto
direzioni: ogni passo consiste nel condurre per un punto di un disegno la parallela a un segmento
dell'altro.

Due avvertenze sulle notazioni. Le forze del sistema si indicano con $\fb_1, \dots, \fb_n$ e le
loro rette d'azione con $a_1, \dots, a_n$: nel piano una forza si può far scorrere lungo la propria
retta d'azione, e il punto di applicazione non interviene. Inoltre nel
\CAPref{cap:statica_CR} il \emph{polo} è il punto rispetto al quale si calcola il momento; qui la
parola indica, secondo l'uso della statica grafica, un punto $P$ del piano delle forze. Dove
occorre il primo significato, si dirà esplicitamente ``il momento rispetto al punto $Q$''.

\section{Il poligono delle forze}
\label{sec:poligono_forze_sgr}

Sia dato un sistema piano di $n$ forze $\fb_1, \dots, \fb_n$, con rette d'azione
$a_1, \dots, a_n$ (\FIGref{fig:costruzione_sgr}\,a); le eventuali coppie si intendono sostituite
ciascuna da due forze opposte su rette parallele. Fissato l'ordine delle forze, nel piano delle
forze le si riporta una di seguito all'altra a partire da un punto $0$: la forza $\fb_i$ è
rappresentata dal segmento orientato che va dal vertice $i-1$ al vertice $i$,
\begin{equation}
	\fb_i = \overrightarrow{(i-1)\,i}, \qquad i = 1, \dots, n .
	\label{eq:lati_poligono_forze_sgr}
\end{equation}
La spezzata $0\,1 \dots n$ è il \emph{poligono delle forze} (\FIGref{fig:costruzione_sgr}\,b), e
il segmento orientato che ne congiunge gli estremi è la risultante \eqref{eq:risultante} del
sistema:
\begin{equation}
	\fb = \sum_{i=1}^{n} \fb_i = \overrightarrow{0\,n} .
	\label{eq:risultante_sgr}
\end{equation}
Dall'ordine in cui le forze sono riportate dipende la forma del poligono, non la risultante.

Il poligono delle forze determina la risultante come vettore libero; non dice dove essa sia
applicata. Nel piano, quando $\fb \neq \zerovb$, il sistema equivale alla sola risultante applicata
sull'asse di sollecitazione, la retta dei punti rispetto ai quali il momento è nullo: per
individuarla occorre tornare al piano di posizione. Per due forze le cui rette d'azione si
incontrano, l'asse di sollecitazione passa per il punto di incontro, nel quale le forze si
trasportano e si compongono. Con più forze si può ripetere l'operazione componendole a due a due,
ma il procedimento cade in difetto quando le rette d'azione sono parallele, e diviene impreciso
quando si incontrano sotto angoli piccoli o fuori dal foglio. Il poligono funicolare risolve il
problema in ogni caso, sostituendo alle composizioni successive altrettante decomposizioni secondo
direzioni scelte a piacere.

\begin{figure}[htbp]
	\centering
	\FiguraSGR{figura-sgr-costruzione}
	\caption{Costruzione del poligono funicolare di un sistema di tre forze: \emph{(a)} piano di
		posizione, con le rette d'azione $a_i$, il poligono funicolare $A_0 \dots A_4$ di lati
		$b_0, \dots, b_3$ e il punto $R$ in cui si incontrano il primo e l'ultimo lato, per il quale
		passa la retta d'azione $a$ della risultante; \emph{(b)} poligono delle forze $0\,1\,2\,3$, con
		il polo $P$ e i raggi polari: il lato $b_i$ è parallelo al raggio $iP$, e la risultante $\fb$
		va da $0$ a $3$.}
	\label{fig:costruzione_sgr}
\end{figure}

\section{Il poligono funicolare}
\label{sec:poligono_funicolare_sgr}

\subsection{La costruzione}
\label{sec:costruzione_sgr}

Nel piano delle forze si sceglie un punto $P$, il \emph{polo}, e lo si congiunge con i vertici del
poligono delle forze: i segmenti $0P, 1P, \dots, nP$ sono i \emph{raggi polari}. Nel piano di
posizione si sceglie un punto $A_0$ e si procede di lato in lato
(\FIGref{fig:costruzione_sgr}):
\begin{itemize}
	\item per $A_0$ si conduce la retta $b_0$ parallela al raggio $0P$, fino a incontrare la retta
	d'azione $a_1$ nel punto $A_1$;
	\item per $A_1$ si conduce la retta $b_1$ parallela al raggio $1P$, fino a incontrare $a_2$ nel
	punto $A_2$;
	\item si prosegue allo stesso modo fino alla retta $b_n$, condotta per $A_n$ parallelamente al
	raggio $nP$, sulla quale si fissa un punto $A_{n+1}$.
\end{itemize}
La spezzata $A_0 A_1 \dots A_{n+1}$ è un \emph{poligono funicolare} del sistema di forze. I punti
$A_1, \dots, A_n$ ne sono i \emph{nodi}, uno su ciascuna retta d'azione, e i segmenti
$b_0, \dots, b_n$ i \emph{lati}: il lato $b_i$ congiunge $A_i$ con $A_{i+1}$ ed è parallelo al
raggio $iP$. Il lato compreso fra le rette d'azione di due forze consecutive è dunque parallelo al
raggio che esce dal vertice comune ai segmenti che le rappresentano.

La costruzione riesce se ogni lato incontra la retta d'azione successiva, cioè se $b_{i-1}$ non è
parallelo ad $a_i$: il polo va scelto fuori dalle rette che contengono i lati del poligono delle
forze. Per il resto la scelta è libera. Fissato l'ordine delle forze, sono arbitrari il polo, cioè
due parametri, e la retta $b_0$ fra le parallele al raggio $0P$, cioè un terzo parametro: i
poligoni funicolari di un sistema di forze sono $\infty^3$.

\subsection{Equivalenza e asse di sollecitazione}
\label{sec:equivalenza_sgr}

Ogni passo della costruzione è una delle operazioni elementari di equivalenza del
\PARref{sec:equivalenza}: traslazione di una forza lungo la retta d'azione, decomposizione,
rimozione di un sistema equilibrato. Trasportata lungo la propria retta d'azione nel nodo $A_i$,
la forza $\fb_i$ si decompone secondo le direzioni dei due lati che vi concorrono. Nel poligono
delle forze la decomposizione è il triangolo di vertici $i-1$, $i$ e $P$:
\begin{equation}
	\fb_i = \overrightarrow{(i-1)\,P} + \overrightarrow{P\,i} ,
	\label{eq:decomposizione_sgr}
\end{equation}
con la prima componente diretta come $b_{i-1}$ e la seconda come $b_i$. Su ciascun lato intermedio
$b_i$, con $1 \le i \le n-1$, vengono così a trovarsi due forze: $\overrightarrow{P\,i}$, che
proviene dalla decomposizione di $\fb_i$ in $A_i$, e $\overrightarrow{i\,P}$, che proviene dalla
decomposizione di $\fb_{i+1}$ in $A_{i+1}$. Hanno la stessa retta d'azione e sono opposte:
formano un sistema equilibrato, che si rimuove. Restano due sole forze,
\begin{equation}
	\overrightarrow{0\,P} \ \text{ sulla retta } b_0 , \qquad
	\overrightarrow{P\,n} \ \text{ sulla retta } b_n ,
	\label{eq:sistema_ridotto_sgr}
\end{equation}
e il sistema dato è equivalente a esse. Con $n$ decomposizioni, tutte eseguite in punti e secondo
direzioni che la scelta del polo rende comodi, le $n$ forze sono ridotte a due.

Se $\fb \neq \zerovb$ e il polo non è allineato con i vertici $0$ e $n$, i lati estremi non sono
paralleli, e le loro rette si incontrano in un punto $R$. Trasportate in $R$ e composte, le due
forze \eqref{eq:sistema_ridotto_sgr} danno
\begin{equation}
	\overrightarrow{0\,P} + \overrightarrow{P\,n} = \overrightarrow{0\,n} = \fb :
	\label{eq:composizione_sgr}
\end{equation}
il sistema equivale alla risultante applicata in $R$, e l'asse di sollecitazione è la retta $a$
condotta per $R$ parallelamente al segmento $0n$. Il punto $R$ cambia con il polo e con il punto
iniziale; la retta $a$ no, perché l'asse di sollecitazione è una proprietà del sistema di forze.
Al variare del poligono funicolare, il punto di incontro del primo e dell'ultimo lato scorre lungo
la retta d'azione della risultante.

\begin{esempio}[Riduzione grafica di un sistema di tre forze]
	\label{ese:riduzione_sgr}
	Si riduca per via grafica il sistema dell'\ESEref{ex:C02_riduzione_piano}: nel piano $(x, y)$,
	$\fb_1 = 2\,\ab_x$ applicata in $(0,\,1)$, $\fb_2 = -\ab_y$ applicata in $(3,\,0)$ e
	$\fb_3 = \ab_x + 2\,\ab_y$ applicata in $(1,\,2)$.

	\begin{svolgimento}
		Riportate le forze una di seguito all'altra a partire dal vertice $0$, il poligono delle forze
		ha i vertici $1$, $2$ e $3$ della \FIGref{fig:esempio_sgr}\,b, e il segmento $0\,3$ fornisce la
		risultante, di componenti $X = 3$ e $Y = 1$. Scelto il polo $P$, per un punto $A_1$ della retta
		$a_1$ si conducono il lato $b_0$, parallelo al raggio $0P$, e il lato $b_1$, parallelo a $1P$,
		fino al nodo $A_2$ su $a_2$; di qui il lato $b_2$, parallelo a $2P$, fino al nodo $A_3$ su
		$a_3$, e infine il lato $b_3$, parallelo a $3P$ (\FIGref{fig:esempio_sgr}\,a). Le rette del
		primo e dell'ultimo lato si incontrano in $R$, e la parallela al segmento $0\,3$ condotta per
		$R$ è l'asse di sollecitazione.

		Sul disegno si legge che l'asse passa per il punto $C = (-0{,}5;\,1{,}5)$, piede della
		perpendicolare condotta dall'origine $O$, a distanza $\overline{OC} \simeq 1{,}58$, e che la
		risultante, applicata su di esso, ruota attorno a $O$ in verso orario: il momento vale
		$\mu_z(O) = -|\fb|\,\overline{OC} = -\sqrt{10}\cdot 5/\sqrt{10} = -5$. Sono i risultati
		dell'esercizio, ottenuti con la riga e la squadra. Cambiando il polo, o il punto $A_1$, cambiano
		il poligono e il punto $R$, ma non la retta $a$.
	\end{svolgimento}
\end{esempio}

\begin{figure}[htbp]
	\centering
	\FiguraSGR{figura-sgr-esempio}
	\caption{Riduzione grafica del sistema dell'\ESEref{ex:C02_riduzione_piano}: \emph{(a)} piano
		di posizione, con le forze nei punti di applicazione, il poligono funicolare e il punto $R$; la
		retta $a$, condotta per $R$ parallelamente alla risultante, è l'asse di sollecitazione e passa
		per il punto $C$; \emph{(b)} poligono delle forze, con il polo e i raggi polari.}
	\label{fig:esempio_sgr}
\end{figure}

\subsection{Forze parallele}
\label{sec:parallele_sgr}

La costruzione non richiede che le rette d'azione delle forze si incontrino, e vale senza
modifiche per le forze parallele, per le quali la composizione a due a due non è possibile. I
vertici del poligono delle forze sono allora allineati su una retta parallela alle forze, la
\emph{linea dei carichi}, e la distanza del polo da essa, letta nella scala delle forze, è la
\emph{distanza polare} $H$: è la componente, uguale per tutti, dei vettori $\overrightarrow{i\,P}$
nella direzione ortogonale alle forze.

Nella \FIGref{fig:parallele_sgr} la costruzione è applicata al caso (b) dell'esempio delle forze
concentrate su un intervallo del \CAPref{cap:statica_CR}
(\FIGref{fig:forze_concentrate_intervallo}\,b): due forze concordi, di intensità $F$ e
$2F$, applicate a un terzo e a due terzi di un intervallo di lunghezza $\ell$. La linea dei carichi
è il segmento $0\,2$, di lunghezza $3F$; il primo e l'ultimo lato si incontrano in $R$, e la
risultante agisce sulla retta $a$ condotta per $R$ parallelamente alle forze, a distanza
$d = 5\ell/9$ dall'origine dell'intervallo: è il baricentro \eqref{eq:baricentro_forze} dei punti
di applicazione, pesati con le intensità.

\begin{figure}[htbp]
	\centering
	\FiguraSGR{figura-sgr-parallele}
	\caption{Risultante di due forze parallele e concordi, di intensità $F$ e $2F$: \emph{(a)} il
		primo e l'ultimo lato del poligono funicolare si incontrano in $R$, sulla retta d'azione $a$
		della risultante, a distanza $d = 5\ell/9$ dall'origine; \emph{(b)} linea dei carichi e polo.}
	\label{fig:parallele_sgr}
\end{figure}

\subsection{La lettura meccanica e i campi}
\label{sec:lettura_meccanica_sgr}

Il nome del poligono viene dalla sua lettura meccanica. Si pensi la spezzata materializzata in una
catena di aste rettilinee incernierate nei nodi, caricata in ciascun nodo $A_i$ dalla forza
$\fb_i$ e trattenuta alle estremità $A \equiv A_0$ e $B \equiv A_{n+1}$ dalle forze
\begin{equation}
	\fb_A = \overrightarrow{P\,0} \ \text{ in } A , \qquad
	\fb_B = \overrightarrow{n\,P} \ \text{ in } B ,
	\label{eq:forze_estremita_sgr}
\end{equation}
dirette secondo il primo e l'ultimo lato (\FIGref{fig:nodi_sgr}\,a). Sono le opposte delle
\eqref{eq:sistema_ridotto_sgr}: le forze date e le forze di estremità formano un sistema
equilibrato, e il poligono delle forze, completato dai due raggi estremi, si chiude nel polo.

Ogni asta, caricata soltanto alle estremità, è un corpo soggetto a due forze, e trasmette una
forza diretta secondo il proprio asse. Il lato $b_i$ esercita sul primo estremo $A_i$ la forza
$\overrightarrow{i\,P}$ e sul secondo estremo $A_{i+1}$ la forza opposta $\overrightarrow{P\,i}$,
di intensità
\begin{equation}
	N_i = \bigl|\overrightarrow{i\,P}\bigr| , \qquad i = 0, \dots, n ,
	\label{eq:forza_lato_sgr}
\end{equation}
pari alla lunghezza del raggio polare, letta nella scala delle forze. Con queste forze ogni nodo è
in equilibrio. Sul nodo $A_i$ agiscono la forza $\fb_i$, la forza $\overrightarrow{i\,P}$
esercitata dal lato $b_i$ e la forza $\overrightarrow{P\,(i-1)}$ esercitata dal lato $b_{i-1}$, e
la loro somma è nulla,
\begin{equation}
	\fb_i + \overrightarrow{i\,P} + \overrightarrow{P\,(i-1)} = \zerovb , \qquad i = 1, \dots, n ,
	\label{eq:equilibrio_nodo_sgr}
\end{equation}
perché è la \eqref{eq:decomposizione_sgr}, letta come chiusura del triangolo di vertici $i-1$, $i$
e $P$. Il poligono delle forze è dunque l'unione di $n$ triangoli, uno per nodo, ciascuno dei
quali ha un raggio in comune con il successivo (\FIGref{fig:nodi_sgr}\,b). Ogni raggio interno è
percorso due volte, in versi opposti, nei due triangoli che lo condividono: sono le forze, uguali
e contrarie, che lo stesso lato esercita sui suoi due nodi di estremità.

\begin{figure}[htbp]
	\centering
	\FiguraSGR{figura-sgr-nodi}
	\caption{Lettura meccanica del poligono funicolare: \emph{(a)} le forze applicate nei nodi e le
		forze di estremità $\fb_A$ e $\fb_B$, dirette secondo il primo e l'ultimo lato; i simboli
		cerchiati sono i campi del piano di posizione, con il campo $P$ dalla parte opposta alle forze;
		\emph{(b)} il poligono delle forze, chiuso nel polo dalle forze di estremità: i triangoli, a
		grigi alterni, sono i poligoni di equilibrio dei nodi $A_1$, $A_2$ e $A_3$, e ogni raggio
		interno è percorso in versi opposti nei due triangoli che lo hanno in comune.}
	\label{fig:nodi_sgr}
\end{figure}

\paragraph{I campi.}
La corrispondenza fra i due disegni si legge a colpo d'occhio con la notazione dei campi, dovuta
a Bow. Nel piano di posizione si riportano le forze nei nodi, tutte dalla stessa parte del
poligono, facendo terminare nel nodo, anziché partirne, quelle che puntano dalla parte opposta. I
lati e le forze dividono allora il piano in campi: uno dalla parte opposta alle forze, al quale si
dà il nome del polo, $P$; gli altri dalla parte delle forze, numerati da $0$ a $n$ in modo che la
forza $\fb_i$ separi il campo $i-1$ dal campo $i$. Il lato $b_i$ separa così il campo $i$ dal
campo $P$, e le forze di estremità, dirette secondo $b_0$ e $b_n$, il campo $P$ dai campi $0$ e
$n$. Con questi nomi:
\begin{itemize}
	\item a ogni campo del piano di posizione corrisponde il vertice del poligono delle forze che
	porta lo stesso nome;
	\item a ogni linea che separa due campi corrisponde il segmento, a essa parallelo, che congiunge
	i due vertici omonimi: alla forza $\fb_i$ il segmento $(i-1)\,i$, al lato $b_i$ il raggio $iP$;
	\item a ogni nodo, attorno al quale si incontrano tre campi, corrisponde il triangolo che ha per
	vertici i tre punti omonimi: al nodo $A_i$ il triangolo $(i-1)\,i\,P$.
\end{itemize}
Punti e regioni si scambiano dunque fra i due disegni, che per questo si dicono \emph{figure
reciproche}. La regola pratica che se ne trae è che il nome di un lato, cioè la coppia dei campi
che esso separa, indica quale segmento del poligono delle forze gli è parallelo e ne misura la
forza.

\begin{oss}[Il verso di lettura]
	\label{oss:verso_lettura_sgr}
	I nomi dei campi fissano anche il verso delle forze. Si giri attorno a un nodo nel verso in
	cui sono stati numerati i campi, che nella \FIGref{fig:nodi_sgr} è quello antiorario: i campi
	che si incontrano, letti nell'ordine, sono i vertici del poligono di equilibrio del nodo,
	percorso nel verso delle forze che agiscono su di esso. Attorno al nodo $A_i$ si incontrano i
	campi $P$, $i-1$ e $i$: il lato $b_{i-1}$ esercita sul nodo la forza
	$\overrightarrow{P\,(i-1)}$, la forza applicata è $\overrightarrow{(i-1)\,i}$, e il lato $b_i$
	esercita la forza $\overrightarrow{i\,P}$. Un lato è teso se la forza che esercita sul nodo se
	ne allontana lungo il lato, cioè se lo tira; è compresso se punta verso il nodo, cioè se lo
	spinge. Per le forze parallele e concordi i lati sono tutti tesi, o tutti compressi, secondo la
	parte della linea dei carichi in cui si è scelto il polo: nel primo caso il poligono è la
	configurazione di una fune caricata dalle forze, nel secondo la linea d'asse di un arco che le
	porta senza inflettersi. Queste due letture, e le costruzioni con cui il poligono si fa
	passare per punti assegnati, sono sviluppate nell'appendice sul poligono funicolare del
	Volume~III.
\end{oss}

\subsection{Il momento del sistema}
\label{sec:momento_sgr}

Il poligono funicolare fornisce anche il momento del sistema rispetto a un punto qualsiasi $Q$ del
piano di posizione (\FIGref{fig:momento_sgr}). Sia $a'$ la retta per $Q$ parallela all'asse di
sollecitazione $a$, e siano $Q'$ e $Q''$ i punti in cui essa incontra le rette del primo e
dell'ultimo lato. I triangoli $R\,Q'Q''$ e $P\,0\,n$ hanno i lati a due a due paralleli, e sono
simili. Nel primo, l'altezza relativa al lato $Q'Q''$ è la distanza fra le rette $a$ e $a'$, cioè
il braccio $d_Q$ della risultante rispetto a $Q$; nel secondo, l'altezza relativa al lato $0n$ è
la distanza $\delta$ del polo dalla retta $0n$, la \emph{distanza polare}. Dalla proporzione
$\overline{Q'Q''} : d_Q = |\fb| : \delta$ segue
\begin{equation}
	|\mu_z(Q)| = |\fb|\,d_Q = \delta\;\overline{Q'Q''} :
	\label{eq:momento_sgr}
\end{equation}
il momento del sistema rispetto a $Q$ è il prodotto della distanza polare, letta nella scala delle
forze, per il segmento che il primo e l'ultimo lato staccano sulla parallela alla risultante
condotta per $Q$, letto nella scala delle lunghezze. Il verso si riconosce dalla posizione di $Q$
rispetto alla retta $a$. Fissato il polo, la distanza polare è la stessa per tutti i punti del
piano, e i segmenti staccati dai lati estremi misurano il momento in un'unica scala: esso cresce
in proporzione alla distanza da $a$, come vuole la legge di variazione
\eqref{eq:variazione_momento}, e si annulla nei punti di $a$, per i quali $Q'$ e $Q''$ coincidono
con $R$.

\begin{figure}[htbp]
	\centering
	\FiguraSGR{figura-sgr-momento}
	\caption{Momento del sistema rispetto a un punto $Q$: \emph{(a)} la retta $a'$, condotta per
		$Q$ parallelamente alla retta d'azione $a$ della risultante, incontra le rette del primo e
		dell'ultimo lato in $Q'$ e $Q''$; \emph{(b)} il triangolo $P\,0\,3$, simile al triangolo
		$R\,Q'Q''$, e la distanza polare $\delta$. Il momento vale $\delta\;\overline{Q'Q''}$.}
	\label{fig:momento_sgr}
\end{figure}

\begin{oss}[Sistemi parziali]
	\label{oss:momento_parziale_sgr}
	La costruzione vale anche per le sole prime $i$ forze, il cui poligono funicolare è la parte
	$A_0 \dots A_{i+1}$ di quello dell'intero sistema. La loro risultante è $\overrightarrow{0\,i}$,
	la sua retta d'azione passa per il punto di incontro delle rette $b_0$ e $b_i$, e il loro
	momento rispetto a $Q$ è il prodotto della distanza del polo dalla retta $0i$ per il segmento
	che $b_0$ e $b_i$ staccano sulla parallela a $0i$ condotta per $Q$. Per un sistema qualsiasi la
	direzione di $0i$ e la distanza del polo da essa cambiano con $i$. Per le forze parallele sono
	le stesse per ogni $i$: la retta $0i$ è la linea dei carichi e la distanza vale $H$. Il
	poligono funicolare, letto lungo le parallele alle forze a partire dalla retta del primo lato, è
	allora il diagramma del momento delle forze che precedono, in scala $H$.
\end{oss}

\section{Le condizioni grafiche di equilibrio}
\label{sec:equilibrio_sgr}

\subsection{Poligoni chiusi}
\label{sec:poligoni_chiusi_sgr}

Sia ora $\fb = \zerovb$. Il poligono delle forze è \emph{chiuso}: il vertice $n$ coincide con il
vertice $0$, e il raggio $nP$ con il raggio $0P$. Il primo e l'ultimo lato del poligono
funicolare, paralleli allo stesso raggio, sono paralleli fra loro
(\FIGref{fig:equilibrio_sgr}\,a), e le due forze \eqref{eq:sistema_ridotto_sgr} alle quali il
sistema si riduce, $\overrightarrow{0\,P}$ sulla retta $b_0$ e
$\overrightarrow{P\,n} = \overrightarrow{P\,0}$ sulla retta $b_n$, sono opposte. Il sistema
equivale a una coppia, di momento
\begin{equation}
	|\mu_z| = N_0\,d ,
	\label{eq:coppia_sgr}
\end{equation}
dove $N_0$ è la lunghezza del raggio $0P$, letta nella scala delle forze, e $d$ la distanza fra le
rette $b_0$ e $b_n$, letta nella scala delle lunghezze. Il momento della coppia è una proprietà
del sistema: se si sposta il polo, $N_0$ e $d$ variano in ragione inversa.

Se anche $d = 0$, l'ultimo lato torna sulla retta del primo: il poligono funicolare è
\emph{chiuso}, la coppia ha braccio nullo e il sistema è equilibrato
(\FIGref{fig:equilibrio_sgr}\,b). Viceversa, un sistema equilibrato ha risultante nulla, quindi
poligono delle forze chiuso, e momento nullo, quindi $d = 0$ per la \eqref{eq:coppia_sgr}.
Dunque \emph{un sistema piano di forze è equilibrato se e solo se sono chiusi tanto il poligono
delle forze quanto il poligono funicolare}. Le due chiusure traducono le due equazioni di
equilibrio \eqref{eq:equilibrio_CR}: la prima l'annullarsi della risultante, la seconda
l'annullarsi del momento. Dei tre casi in cui si classificano i sistemi piani, la forza unica ha
il poligono delle forze aperto, la coppia ha il poligono delle forze chiuso e il poligono
funicolare aperto, il sistema nullo li ha chiusi entrambi.

Nella lettura meccanica il poligono funicolare chiuso è un anello di aste in equilibrio sotto le
sole forze applicate nei nodi, senza forze di estremità.

\begin{figure}[htbp]
	\centering
	\FiguraSGR{figura-sgr-equilibrio}
	\caption{Sistema di tre forze a risultante nulla: \emph{(a)} la retta d'azione di $\fb_3$ non
		passa per il punto $K$ in cui concorrono le altre due: il primo e l'ultimo lato del poligono
		funicolare sono paralleli, a distanza $d$, e il sistema equivale a una coppia; \emph{(b)} le tre
		rette d'azione concorrono in $K$: l'ultimo lato torna sulla retta del primo, il poligono
		funicolare si chiude e il sistema è equilibrato; \emph{(c)} il poligono delle forze, chiuso,
		comune ai due casi.}
	\label{fig:equilibrio_sgr}
\end{figure}

\subsection{Due e tre forze}
\label{sec:tre_forze_sgr}

Per i sistemi di due e di tre forze le condizioni di chiusura si riducono a quelle già enunciate
nel \PARref{sec:equivalenza}, e non richiedono il poligono funicolare.
\begin{itemize}
	\item \emph{Due forze} sono in equilibrio se e solo se sono opposte e hanno la stessa retta
	d'azione: il poligono delle forze è un segmento percorso due volte. Ne discende che un corpo
	scarico, vincolato al resto della struttura da due sole cerniere, trasmette una forza diretta
	secondo la congiungente delle cerniere: è la \emph{biella}.
	\item \emph{Tre forze} non parallele sono in equilibrio se e solo se il loro poligono delle forze
	è un triangolo chiuso e le loro rette d'azione concorrono in un punto
	(\FIGref{fig:equilibrio_sgr}). Se infatti due di esse si incontrano in $K$, la loro risultante
	passa per $K$, e la terza, che deve esserle opposta sulla stessa retta, passa anch'essa per
	$K$.
\end{itemize}
Queste due proposizioni risolvono, senza altro, gran parte dei problemi di equilibrio dei sistemi
isostatici: una retta d'azione incognita si determina perché deve passare per il punto in cui
concorrono le altre due, e le intensità si leggono sul triangolo delle forze.

\section{Le reazioni vincolari}
\label{sec:reazioni_sgr}

Le condizioni di chiusura servono a determinare le reazioni dei vincoli. Di ciascuna reazione il
\CAPref{cap:vincoli_CR} dice che cosa è noto a priori: la reazione di un carrello è diretta
secondo l'asse del carrello, e ne è incognita la sola intensità \eqref{eq:reazione_carrello}; la
reazione di una cerniera è una forza che passa per la cerniera, di direzione e intensità incognite
\eqref{eq:reazione_cerniera}. Forze attive e reazioni formano un sistema equilibrato, e i loro
due poligoni devono chiudersi: da questa condizione si ricavano, in un corpo isostatico, le tre
incognite.

\subsection{Forze parallele: il lato di chiusura}
\label{sec:chiusura_sgr}

Il caso più semplice è quello di un corpo soggetto a forze parallele $\fb_1, \dots, \fb_n$ e
vincolato in due punti $A$ e $B$ in modo che anche le reazioni $\rb_A$ e $\rb_B$ siano parallele
alle forze: la trave appoggiata soggetta a carichi verticali (\FIGref{fig:chiusura_sgr}). Delle
reazioni sono note le rette d'azione $a_A$ e $a_B$ e incognite le intensità. Forze e reazioni,
prese nell'ordine $\fb_1, \dots, \fb_n, \rb_B, \rb_A$, formano un sistema equilibrato.

Si costruisca, con un polo qualsiasi, il poligono funicolare delle forze date, prendendone gli
estremi sulle rette d'azione delle reazioni: $A_0$ su $a_A$ e $A_{n+1}$ su $a_B$. I due estremi
sono i nodi delle reazioni, e per chiudere il poligono funicolare dell'intero sistema manca un
solo lato, quello che da $A_{n+1}$ torna ad $A_0$: è il segmento $c = A_0A_{n+1}$, il \emph{lato
di chiusura}. Nel poligono delle forze gli corrisponde un raggio a esso parallelo. Condotta per
il polo la parallela a $c$, essa incontra la linea dei carichi in un punto $m$, che è il vertice
mancante del poligono chiuso $0\,1 \dots n\,m\,0$:
\begin{equation}
	\rb_B = \overrightarrow{n\,m} , \qquad \rb_A = \overrightarrow{m\,0} .
	\label{eq:reazioni_sgr}
\end{equation}
\emph{Il raggio parallelo al lato di chiusura divide la linea dei carichi nelle due reazioni.}

Le reazioni sono determinate dall'equilibrio, e il punto $m$ non dipende dunque dal polo né dal
punto iniziale: al variare del poligono funicolare il lato di chiusura cambia, ma il raggio a
esso parallelo passa sempre per lo stesso punto della linea dei carichi.

Con il lato di chiusura si completano anche i campi del piano di posizione
(\PARref{sec:lettura_meccanica_sgr}). Riportate nei nodi le forze e le reazioni, tutte all'esterno
del poligono chiuso, questo racchiude il campo $P$. Dalla parte dei lati $b_0, \dots, b_n$ stanno i
campi $0, \dots, n$, separati dalle forze; dalla parte del lato di chiusura, fra le due reazioni,
sta un nuovo campo, al quale si dà il nome $m$ (\FIGref{fig:chiusura_sgr}\,a). Il lato di chiusura
separa il campo $m$ dal campo $P$, la reazione $\rb_A$ il campo $m$ dal campo $0$, la reazione
$\rb_B$ il campo $n$ dal campo $m$: nel poligono delle forze sono il raggio $mP$ e i segmenti
$m\,0$ e $n\,m$. Attorno al nodo $A_0$ si incontrano i campi $m$, $0$ e $P$, e il triangolo
$m\,0\,P$ è il suo poligono di equilibrio: nella lettura meccanica il lato di chiusura è un'asta
che collega gli estremi della catena, e il raggio $mP$ ne misura la forza.

\begin{figure}[htbp]
	\centering
	\FiguraSGR{figura-sgr-chiusura}
	\caption{Reazioni di una trave appoggiata soggetta a forze parallele: \emph{(a)} il poligono
		funicolare delle forze, con gli estremi $A_0$ e $A_4$ sulle rette d'azione $a_A$ e $a_B$ delle
		reazioni, e il lato di chiusura $c$; $\eta$ è la distanza del poligono dal lato di chiusura,
		misurata nella direzione delle forze. Forze e reazioni sono riportate anche nei nodi del
		poligono chiuso; i simboli cerchiati sono i campi, con $P$ all'interno del poligono e $m$
		dalla parte del lato di chiusura; \emph{(b)} il raggio condotto per il polo parallelamente a
		$c$ incontra la linea dei carichi in $m$ e la divide nelle reazioni
		$\rb_A = \protect\overrightarrow{m\,0}$ e $\rb_B = \protect\overrightarrow{3\,m}$.}
	\label{fig:chiusura_sgr}
\end{figure}

\begin{esempio}[Trave con sbalzo]
	\label{ese:sbalzo_sgr}
	La campata $ABE$ della trave di Gerber del \PARref{subsec:gerber} è vincolata da una cerniera
	in $A$ e da un carrello in $B$, a distanza $\ell$, e porta la forza $p\ell/2$, diretta verso il
	basso, all'estremo $E$ di uno sbalzo di lunghezza $\ell/2$. Se ne determinino graficamente le
	reazioni.

	\begin{svolgimento}
		La retta d'azione $a_1$ della forza è esterna alle rette $a_A$ e $a_B$ delle reazioni
		(\FIGref{fig:sbalzo_sgr}\,a). Il poligono funicolare ha due soli lati: $b_0$, condotto dal punto
		$A_0$ di $a_A$ parallelamente al raggio $0P$ fino al nodo $A_1$ su $a_1$, e $b_1$, parallelo
		al raggio $1P$, che per incontrare $a_B$ va prolungato da $A_1$ all'indietro, fino al punto
		$A_2$. Il lato di chiusura è $c = A_0A_2$. Il raggio condotto per il polo parallelamente a $c$
		incontra la linea dei carichi nel punto $m$, che cade \emph{fuori} dal segmento $0\,1$, dalla
		parte del vertice $0$ (\FIGref{fig:sbalzo_sgr}\,b).

		La reazione $\rb_B = \overrightarrow{1\,m}$ è diretta verso l'alto ed è lunga una volta e
		mezza il segmento $0\,1$: $V_B = \tfrac{3}{4}\,p\ell$. La reazione
		$\rb_A = \overrightarrow{m\,0}$ è diretta verso il basso ed è lunga la metà:
		$V_A = -\tfrac{1}{4}\,p\ell$. Sono i valori ottenuti nel capitolo con le equazioni cardinali.
		Quando il punto $m$ cade fuori dalla linea dei carichi, una delle reazioni è concorde con le
		forze: il vincolo in $A$ non sostiene la trave, la trattiene.
	\end{svolgimento}
\end{esempio}

\begin{figure}[htbp]
	\centering
	\FiguraSGR{figura-sgr-sbalzo}
	\caption{Reazioni di una trave con sbalzo: \emph{(a)} il poligono funicolare ha i lati $b_0$ e
		$b_1$, e l'ultimo incontra la retta $a_B$ nel suo prolungamento all'indietro; $c$ è il lato di
		chiusura; \emph{(b)} il raggio parallelo a $c$ incontra la linea dei carichi in $m$, fuori dal
		segmento $0\,1$: la reazione $\rb_B = \protect\overrightarrow{1\,m}$ è rivolta verso l'alto,
		la reazione $\rb_A = \protect\overrightarrow{m\,0}$ verso il basso.}
	\label{fig:sbalzo_sgr}
\end{figure}

\begin{oss}[Il poligono funicolare come diagramma dei momenti]
	\label{oss:momento_flettente_sgr}
	Si applichi l'\OSSref{oss:momento_parziale_sgr} al sistema equilibrato formato dalle forze e
	dalle reazioni, ordinato a partire da $\rb_A$: il suo poligono funicolare comincia con la retta
	del lato di chiusura e prosegue con i lati $b_0, \dots, b_n$. Per una sezione della trave
	compresa fra le rette $a_i$ e $a_{i+1}$, il momento delle forze che la precedono, reazione
	compresa, è il prodotto della distanza polare per il segmento $\eta$ che il lato di chiusura e
	il lato $b_i$ staccano sulla parallela alle forze condotta per la sezione:
	\begin{equation}
		H\,\eta = |M| .
		\label{eq:momento_flettente_sgr}
	\end{equation}
	Ma il momento delle forze che precedono una sezione è, in valore assoluto, il momento flettente
	della trave (\CAPref{cap:Diagrammi_CS}): \emph{il poligono funicolare, misurato dal lato di
	chiusura nella direzione delle forze, è il diagramma del momento flettente della trave
	appoggiata, in scala $H$}. Nella \FIGref{fig:chiusura_sgr} il momento ha lo stesso segno lungo
	tutta la trave; nella \FIGref{fig:sbalzo_sgr} il poligono chiuso è il triangolo $A_0A_1A_2$, e
	il momento, nullo in $A$ e in $E$, è massimo in valore assoluto sull'appoggio $B$.
\end{oss}

\subsection{Forze qualsiasi: cerniera e carrello}
\label{sec:cerniera_carrello_sgr}

Sia ora il corpo soggetto a forze qualsiasi e vincolato con una cerniera in $A$ e un carrello in
$B$. Della reazione $\rb_B$ è nota la retta d'azione $a_B$, l'asse del carrello; della reazione
$\rb_A$ è noto soltanto un punto, la cerniera. Il lato di chiusura si costruisce ancora, con una
sola avvertenza: si conduce il primo lato del poligono funicolare per $A$, l'unico punto noto
della retta d'azione di $\rb_A$, e si prende l'estremo $A_{n+1}$ sulla retta $a_B$, sicché
$c = A\,A_{n+1}$. Nel poligono delle forze il punto $m$ è l'intersezione della parallela a $c$
condotta per il polo con la parallela ad $a_B$ condotta per il vertice $n$, e valgono ancora le
\eqref{eq:reazioni_sgr}.

Quando le forze attive si riducono a una sola forza, o sono state composte nella loro risultante,
basta la proposizione sulle tre forze (\PARref{sec:tre_forze_sgr}): la retta d'azione della
forza attiva e l'asse del carrello si incontrano in un punto $K$, e la reazione della cerniera,
che deve passare per $K$, ha la retta $AK$. Il triangolo delle forze, costruito conducendo per
gli estremi della forza attiva le parallele ad $a_B$ e ad $AK$, fornisce le due reazioni.

\begin{esempio}[Tre forze concorrenti]
	\label{ese:tre_forze_sgr}
	Si determinino graficamente le reazioni esterne del sistema dell'\ESEref{ex:anello_vincolato_TA}.

	\begin{svolgimento}
		Per le reazioni esterne il sistema si comporta come un solo corpo rigido, vincolato a terra
		dalla cerniera in $B$ e dalla biella $AC$, che è un carrello con l'asse diretto secondo la
		propria corda (\FIGref{fig:tre_forze_sgr}\,a). I carichi sono il carico $p$, riferito alla
		proiezione orizzontale del tratto $ED$, di risultante $p\ell$ sulla verticale a distanza $\ell$
		da $A$, e la coppia oraria $p\ell^2$ applicata in $G$.

		Una forza e una coppia si compongono in una sola forza, uguale alla prima e spostata
		parallelamente a sé stessa della distanza $|\mu|/|\fb| = \ell$, dalla parte per la quale il
		momento della forza spostata rispetto ai punti della retta originaria ha il verso della coppia:
		è l'operazione inversa del trasporto di una forza. La forza $p\ell$ agisce dunque sulla
		verticale $a$ a distanza $2\ell$ da $A$. Questa retta incontra l'asse $a_A$ della biella nel
		punto $K$, e la reazione della cerniera ha la retta $a_B = BK$.

		Nel triangolo delle forze (\FIGref{fig:tre_forze_sgr}\,b) il segmento $0\,1$ rappresenta la
		forza $p\ell$; la parallela ad $a_B$ per il vertice $1$ e la parallela ad $a_A$ per il vertice
		$0$ si incontrano in $m$, e $\rb_B = \overrightarrow{1\,m}$, $\rb_A = \overrightarrow{m\,0}$.
		Proiettando i due lati sugli assi si leggono le componenti: $V_B = \tfrac{4}{5}\,p\ell$ e
		$H_B = \tfrac{1}{10}\,p\ell$ verso sinistra; $V_A = \tfrac{1}{5}\,p\ell$ e
		$H_A = \tfrac{1}{10}\,p\ell$ verso destra. Sono i valori ottenuti nell'esercizio.
	\end{svolgimento}
\end{esempio}

\begin{figure}[htbp]
	\centering
	\FiguraSGR{figura-sgr-tre-forze}
	\caption{Reazioni del sistema dell'\ESEref{ex:anello_vincolato_TA}: \emph{(a)} la risultante
		$p\ell$ del carico distribuito e la coppia $p\ell^2$ equivalgono alla forza $p\ell$ sulla retta
		$a$, spostata di $\ell$; la retta $a$ e l'asse $a_A$ della biella $AC$ si incontrano in $K$,
		per il quale passa la retta d'azione $a_B$ della reazione in $B$; \emph{(b)} triangolo delle
		forze.}
	\label{fig:tre_forze_sgr}
\end{figure}

\subsection{Sistemi a tre cerniere}
\label{sec:tre_cerniere_sgr}

In un sistema a tre cerniere le reazioni esterne sono quattro, e la sola chiusura dei poligoni
dell'intero sistema non basta. La proposizione sulle due forze fornisce la condizione mancante
quando uno dei due corpi è scarico: esso è una biella, e la reazione della sua cerniera a terra è
diretta secondo la congiungente con la cerniera interna. Il problema torna allora a quello delle
tre forze.

Nel portale della \FIGref{fig:tre_cerniere_sgr}, che ha la geometria di quello del
\PARref{sec:arco_tre_cerniere}, agisce una sola forza $F$ sul traverso di sinistra. Il corpo
$DEB$ è scarico, e la reazione in $B$ ha la retta $a_B = BD$. Questa incontra la retta d'azione
della forza in $K$, e la reazione in $A$ ha la retta $a_A = AK$. Il triangolo delle forze dà
$\rb_B$, di componenti $F/4$ verso l'alto e $F/4$ verso sinistra, e $\rb_A$, di componenti
$3F/4$ verso l'alto e $F/4$ verso destra: lo stesso risultato si ottiene dall'equilibrio alla
rotazione dell'intero portale attorno ad $A$ e del corpo $DEB$ attorno a $D$.

Se sono carichi entrambi i corpi, si considerano separatamente i carichi dell'uno e dell'altro:
in ciascuno dei due problemi un corpo è scarico, e le reazioni cercate sono la somma di quelle
dei due problemi, che si esegue riportando i due triangoli nello stesso piano delle forze. La
simmetria, quando c'è, abbrevia il procedimento: in un sistema simmetrico e simmetricamente
caricato la forza trasmessa dalla cerniera interna è ortogonale all'asse di simmetria, e ciascuna
metà è soggetta a tre forze, di cui sono note due rette d'azione.

\begin{figure}[htbp]
	\centering
	\FiguraSGR{figura-sgr-tre-cerniere}
	\caption{Portale a tre cerniere con una forza sul traverso di sinistra: \emph{(a)} il corpo
		$DEB$, scarico, trasmette una forza diretta secondo $BD$; le rette d'azione della forza e
		delle due reazioni concorrono in $K$; \emph{(b)} triangolo delle forze, con
		$\rb_B = \protect\overrightarrow{1\,m}$ e $\rb_A = \protect\overrightarrow{m\,0}$.}
	\label{fig:tre_cerniere_sgr}
\end{figure}

\section{I sistemi reticolari}
\label{sec:reticolari_sgr}

Il poligono funicolare, letto come catena di aste, è il più semplice dei sistemi reticolari, e i
triangoli dei suoi nodi sono equilibri di nodi nei quali concorrono due aste. La notazione dei
campi è nata per le travature reticolari, e vi rende il metodo dei nodi (\PARref{sec:metodo_nodi})
una costruzione grafica.

\subsection{Il nodo canonico}
\label{sec:nodo_canonico_sgr}

Si consideri una travatura reticolare caricata nei nodi e liberata dai vincoli, con le reazioni
esterne già determinate. Si fissa un verso di percorrenza, per esempio quello orario, e si dà un
nome ai campi:
\begin{itemize}
	\item i campi \emph{esterni}, compresi fra due forze esterne consecutive, si numerano nell'ordine
	in cui si incontrano girando attorno alla travatura nel verso fissato, e nello stesso ordine si
	riportano le forze esterne nel poligono delle forze, che è chiuso;
	\item i campi \emph{interni} sono le maglie, e si indicano con lettere minuscole,
	$a, b, \dots$, per distinguerli dai primi.
\end{itemize}
Ogni asta separa due campi, e ogni nodo è circondato dai campi che le aste e le forze concorrenti
in esso delimitano. Come nel poligono funicolare, alla forza di un'asta corrisponde, nel piano
delle forze, il segmento parallelo all'asta che congiunge i vertici omonimi dei due campi; e al
nodo corrisponde il poligono chiuso che ha per vertici i campi che lo circondano, letti nel verso
fissato (\OSSref{oss:verso_lettura_sgr}).

In un \emph{nodo canonico} sono incognite le forze di due sole aste. I campi adiacenti alle forze
note hanno vertici già segnati nel piano delle forze; le due aste incognite delimitano un campo
il cui vertice manca. Lo si trova all'intersezione delle parallele alle due aste, condotte per i
vertici dei campi, già noti, che esse separano da quello. Nella
\FIGref{fig:nodo_canonico_sgr} il procedimento è applicato al nodo $D$ della capriata
dell'\ESEref{ese:capriata_palladiana}: attorno al nodo, in verso orario, si incontrano i campi
$0$, $1$, $b$ e $a$. Sono noti il segmento $0\,1$, che è la forza applicata, e
il segmento $a\,0$, che è la forza dell'asta $DA$, determinata nel nodo precedente. Il
vertice $b$ è l'intersezione della parallela a $DF$ per il vertice $1$ con la parallela
a $DB$ per il vertice $a$.

Il poligono $0\,1\,b\,a$, percorso in questo ordine, dà le forze che agiscono
sul nodo, con il loro verso. L'asta $DF$ esercita la forza $\overrightarrow{1\,b}$, che
punta verso il nodo: l'asta lo spinge, ed è compressa. Lo stesso vale per l'asta $DB$, con la
forza $\overrightarrow{b\,a}$. Un'asta è tesa quando la forza che esercita sul
nodo se ne allontana lungo l'asta.

\begin{figure}[htbp]
	\centering
	\FiguraSGR{figura-sgr-nodo-canonico}
	\caption{Equilibrio grafico di un nodo canonico, il nodo $D$ della capriata della
		\FIGref{fig:cremoniano_sgr}: \emph{(a)} il nodo con i campi che lo circondano e le forze che
		agiscono su di esso; la freccia curva indica il verso di lettura; \emph{(b)} il poligono di
		equilibrio: sono noti i lati $a\,0$ e $0\,1$, e il vertice $b$ è
		l'intersezione delle parallele alle due aste incognite.}
	\label{fig:nodo_canonico_sgr}
\end{figure}

\subsection{Il diagramma cremoniano}
\label{sec:cremoniano_sgr}

Procedendo di nodo canonico in nodo canonico, nell'ordine della propagazione a cascata del metodo
dei nodi, ogni nodo aggiunge un vertice al piano delle forze. I poligoni dei nodi non si disegnano
separati: hanno i lati in comune, perché ogni asta separa gli stessi due campi vista dall'uno e
dall'altro dei suoi nodi, e si compongono in una sola figura, il \emph{diagramma cremoniano},
nella quale ogni forza d'asta compare una volta sola. Il poligono delle forze con il polo e i
raggi polari ne è il caso particolare relativo al poligono funicolare.

\begin{esempio}[Il diagramma cremoniano della capriata palladiana]
	\label{ese:cremoniano_sgr}
	Si determinino graficamente gli sforzi normali nelle aste della capriata
	dell'\ESEref{ese:capriata_palladiana}, soggetta a tre forze verticali di intensità $F$ nei nodi
	$D$, $F$ ed $E$.
	
	\begin{svolgimento}
		Le reazioni sono verticali e, per simmetria, valgono $\tfrac{3}{2}F$. Girando attorno alla
		travatura in verso orario a partire dalla reazione in $A$, le forze esterne si incontrano
		nell'ordine $\rb_A$, $\fb_1$, $\fb_2$, $\fb_3$, $\rb_C$, e delimitano i campi esterni $0$, $1$,
		$2$, $3$ e $m$; le quattro maglie sono i campi interni $a$, $b$,
		$c$ e $d$ (\FIGref{fig:cremoniano_sgr}\,a). Il poligono delle forze esterne
		è la linea dei carichi $0\,1\,2\,3$, percorsa verso il basso dalle forze e risalita dalle
		reazioni $\rb_C = \overrightarrow{3\,m}$ e $\rb_A = \overrightarrow{m\,0}$, con $m$ nel punto
		medio.
		
		I vertici dei campi interni si determinano nell'ordine dei nodi canonici.
		\begin{itemize}
			\item \emph{Nodo $A$} (campi $m$, $0$, $a$): il vertice $a$ è l'intersezione
			della parallela ad $AD$ per $0$ con la parallela ad $AB$ per $m$.
			\item \emph{Nodo $D$} (campi $0$, $1$, $b$, $a$): il vertice $b$ è
			l'intersezione della parallela a $DF$ per $1$ con la parallela a $DB$ per $a$.
			\item \emph{Nodo $F$} (campi $1$, $2$, $c$, $b$): il vertice
			$c$ è l'intersezione della parallela a $FE$ per $2$ con la parallela a $FB$ per
			$b$.
			\item \emph{Nodo $E$} (campi $2$, $3$, $d$, $c$): il vertice
			$d$ è l'intersezione della parallela a $EC$ per $3$ con la parallela a $EB$ per
			$c$; per la simmetria coincide con $a$.
		\end{itemize}
		I nodi $C$ e $B$ non aggiungono vertici: i loro poligoni, $3\,m\,d$ e
		$m\,a\,b\,c\,d$, sono già tracciati e devono
		chiudersi, e questa è la verifica del disegno.
		
		A ciascuna asta corrisponde il segmento che congiunge i vertici dei due campi che essa separa
		(\FIGref{fig:cremoniano_sgr}\,b):
		\begin{equation*}
			\begin{array}{llll}
				AD:\ 0\,a , & DF:\ 1\,b , & FE:\ 2\,c , & CE:\ 3\,d , \\[2pt]
				DB:\ a\,b , & BF:\ b\,c , &
				EB:\ c\,d , & AB:\ a\,m , \quad BC:\ d\,m .
			\end{array}
		\end{equation*}
		Le lunghezze, lette nella scala delle forze, danno le intensità: $\tfrac{3\sqrt{2}}{2}F$ per
		le falde $AD$ e $CE$, $\sqrt{2}\,F$ per $DF$ e $FE$, $\tfrac{\sqrt{2}}{2}F$ per le saette $DB$
		ed $EB$, $F$ per il monaco $BF$ e $\tfrac{3}{2}F$ per la catena $AB$, $BC$. Il segno si legge
		girando attorno a un nodo. Attorno ad $A$, in verso orario, i campi sono $m$, $0$,
		$a$: l'asta $AD$ esercita sul nodo la forza $\overrightarrow{0\,a}$, che punta
		verso $A$, ed è compressa; l'asta $AB$ esercita la forza $\overrightarrow{a\,m}$, che
		si allontana da $A$ lungo l'asta, ed è tesa. Risultano compresse le falde e le saette, tesi la
		catena e il monaco: sono i valori e i segni dell'esempio.
	\end{svolgimento}
\end{esempio}

\begin{figure}[htbp]
	\centering
	\FiguraSGR{figura-sgr-cremoniano}
	\caption{Diagramma cremoniano della capriata dell'\ESEref{ese:capriata_palladiana}:
		\emph{(a)} la travatura liberata dai vincoli, con le forze $\fb_1$, $\fb_2$, $\fb_3$, di
		intensità $F$, le reazioni e i campi: quelli esterni sono numerati in verso orario, quelli
		interni, le maglie, sono indicati con lettere; \emph{(b)} il diagramma: a ogni asta
		corrisponde il segmento che
		congiunge i vertici dei due campi che essa separa; tratto spesso per le aste compresse, sottile
		per le tese. I vertici $a$ e $d$ coincidono per la simmetria.}
	\label{fig:cremoniano_sgr}
\end{figure}

\begin{oss}[Aste scariche e limiti del procedimento]
	\label{oss:limiti_cremoniano_sgr}
	Un'asta è scarica quando i due campi che essa separa hanno lo stesso vertice nel piano delle
	forze, e il segmento che le corrisponde si riduce a un punto: i criteri di riconoscimento delle
	aste scariche si leggono così sul diagramma. Il procedimento richiede, a ogni passo, un nodo
	canonico. Dove la propagazione a cascata si arresta, come nelle travature che non si risolvono
	con il solo metodo dei nodi, si interrompe anche il diagramma, e la forza di un'asta va prima
	determinata con una sezione di Ritter (\PARref{sec:metodo_ritter}).
\end{oss}

\section{Il parallelo con la cinematica}
\label{sec:dualita_sgr}

Il volume insiste sulla corrispondenza fra la cinematica e la statica del corpo rigido. Nel piano
essa prende la forma più semplice: a una forza corrisponde una rotazione, alla retta d'azione il
centro di rotazione, alla coppia, forza di retta d'azione impropria, la traslazione, rotazione di
centro improprio; il momento rispetto a un punto e lo spostamento di un punto obbediscono alla
stessa legge di variazione. Le costruzioni di questa appendice hanno tutte, in questa
corrispondenza, la loro controparte nella cinematica grafica del
\CAPref{cap:cinematica_infinitesima} (\TABref{tab:dualita_sgr}).

\begin{table}[htbp]
	\centering
	\caption{Corrispondenza fra le costruzioni della statica grafica e quelle della cinematica
		grafica nel piano.}
	\label{tab:dualita_sgr}
	\begin{tabular}{p{0.46\textwidth}p{0.46\textwidth}}
		\hline
		\emph{Statica} & \emph{Cinematica} \\
		\hline
		forza, di intensità $F$ & rotazione, di ampiezza $\vartheta$ \\
		retta d'azione & centro di rotazione \\
		coppia (retta d'azione impropria) & traslazione (centro improprio) \\
		momento rispetto a un punto & spostamento di un punto \\
		risultante, somma delle forze & rotazione composta, somma delle rotazioni \\
		lati del poligono funicolare & corpi di una catena \\
		nodo fra due lati, sulla retta d'azione di una forza & centro relativo di due corpi \\
		incontro del primo e dell'ultimo lato, sulla retta d'azione della risultante &
		centro assoluto dell'ultimo corpo \\
		tre forze in equilibrio: tre rette per un punto & tre corpi: tre centri relativi su una retta \\
		\hline
	\end{tabular}
\end{table}

\subsection{Poligono funicolare e diagramma degli spostamenti}
\label{sec:spostamenti_sgr}

Si riprendano le due forze parallele della \FIGref{fig:parallele_sgr}, e si scelga il polo sulla
retta condotta per il vertice $0$ ortogonalmente alle forze, a distanza $H$ dalla linea dei
carichi (\FIGref{fig:dualita_sgr}\,a). Il primo lato $b_0$ è ortogonale alle forze; passando per
il nodo $A_1$ il poligono devia di $F/H$, e passando per $A_2$ di $2F/H$. Per
l'\OSSref{oss:momento_parziale_sgr} la distanza del lato $b_i$ dalla retta del primo lato,
misurata nella direzione delle forze in un punto $Q$, è il momento rispetto a $Q$ delle prime $i$
forze, diviso per $H$. Il lato $b_1$ incontra la retta di $b_0$ in $A_1$, sulla retta d'azione di
$\fb_1$, dove il momento di questa è nullo; il lato $b_2$ la incontra in $R$, sulla retta
d'azione della risultante, dove è nullo il momento delle due forze.

Si consideri ora una catena di tre corpi (\FIGref{fig:dualita_sgr}\,b): il corpo $0$ è fermo, il
corpo $1$ ruota di $\vartheta_1$ attorno al punto $C_1$, e il corpo $2$ ruota rispetto al corpo
$1$ attorno al punto $C_{12}$, il loro centro relativo (\PARref{sec:CR_relativo}). Il diagramma
degli spostamenti trasversali alla retta dei centri, già impiegato nella
\FIGref{fig:composizione_2_rot}, è una spezzata che devia di $\vartheta_1$ in $C_1$ e della
rotazione relativa $\vartheta_2 - \vartheta_1$ in $C_{12}$: lo spostamento di un punto del corpo
$i$ è la somma dei prodotti delle rotazioni che precedono per le distanze dai rispettivi centri.
La retta del corpo $1$ incontra la configurazione iniziale in $C_1$, il suo centro assoluto; la
retta del corpo $2$ la incontra in $C_2$, il centro assoluto del corpo $2$, che è il centro della
rotazione composta \eqref{eq:comp_rot_baricentro}.

Se le rotazioni stanno fra loro come le forze, le due spezzate sono, a meno di un ribaltamento
attorno all'orizzontale, la stessa figura. Il poligono
funicolare delle forze parallele è il diagramma dei loro momenti, come la spezzata delle
rotazioni è il diagramma degli spostamenti; ai lati corrispondono i corpi, ai nodi i centri
relativi, al punto $R$ il centro assoluto dell'ultimo corpo. L'allineamento dei centri $C_1$,
$C_{12}$ e $C_2$ e il rapporto in cui il centro relativo divide il segmento dei due centri
assoluti sono, letti sulla \FIGref{fig:dualita_sgr}\,a, la regola con cui la risultante di due
forze parallele divide la distanza fra le loro rette d'azione in parti inversamente proporzionali
alle intensità.

\begin{figure}[htbp]
	\centering
	\FiguraSGR{figura-sgr-dualita}
	\caption{Poligono funicolare e diagramma degli spostamenti: \emph{(a)} le forze $F$ e $2F$
		della \FIGref{fig:parallele_sgr}, con il polo scelto in modo che il primo lato sia ortogonale
		alle forze: il lato $b_2$ incontra la retta di $b_0$ in $R$, sulla retta d'azione della
		risultante; \emph{(b)} catena di tre corpi, con il corpo $0$ fermo: il corpo $1$ ruota di
		$\vartheta_1$ attorno a $C_1$, il corpo $2$ ruota rispetto al corpo $1$ di $2\vartheta_1$
		attorno al centro relativo $C_{12}$; la retta del corpo $2$ incontra la configurazione
		iniziale nel suo centro assoluto $C_2$. Le due spezzate hanno, a meno del verso, le stesse
		pendenze, e i punti
		$A_1$, $A_2$ e $R$ stanno sulle verticali di $C_1$, $C_{12}$ e $C_2$.}
	\label{fig:dualita_sgr}
\end{figure}

\subsection{Rette concorrenti e centri allineati}
\label{sec:concorrenza_sgr}

Nella corrispondenza le rette si scambiano con i punti, e il concorso di più rette in un punto
con l'allineamento di più punti su una retta.
\begin{itemize}
	\item Tre forze in equilibrio hanno rette d'azione concorrenti (\PARref{sec:tre_forze_sgr}); i
	tre centri di rotazione relativi di tre corpi sono allineati
	(\PARref{sec:allineamento_relativi}). In entrambi i casi la proposizione serve a determinare
	l'elemento incognito, una retta o un punto, dagli altri due.
	\item Per un carrello, la retta condotta per il punto di applicazione nella direzione dell'asse è
	a un tempo il luogo dei centri di rotazione che il vincolo consente e la retta d'azione della
	reazione che esso esplica; per una cerniera, il punto fisso è il centro di rotazione, e per esso
	passa la reazione (\CAPref{cap:vincoli_CR}).
	\item Un corpo vincolato da tre carrelli i cui assi concorrono in un punto è labile, perché quel
	punto è un centro di rotazione consentito da tutti e tre. Per la stessa ragione le tre reazioni,
	concorrenti, non possono equilibrare una forza che non passi per quel punto: la condizione
	grafica di labilità e la condizione grafica di impossibilità dell'equilibrio sono la stessa
	figura.
\end{itemize}
La radice comune è il principio dei lavori virtuali, per il quale la reazione di un vincolo non
compie lavoro negli spostamenti che il vincolo consente (\CAPref{cap:dualismo_CR}): una forza non
compie lavoro in una rotazione se e solo se la sua retta d'azione passa per il centro.

\begin{oss}[Nota storica]
	\label{oss:nota_storica_sgr}
	La regola del triangolo delle forze risale a Simon Stevin, alla fine del Cinquecento; il
	poligono delle forze e il poligono funicolare compaiono insieme nella \emph{Nouvelle mécanique}
	di Pierre Varignon, pubblicata postuma nel 1725. La statica grafica come disciplina sistematica
	nasce con \emph{Die graphische Statik} di Karl Culmann, del 1866. La reciprocità fra il piano di
	posizione e il piano delle forze fu riconosciuta da James Clerk Maxwell nel 1864 e sviluppata
	da Luigi Cremona nel 1872, dal quale prende nome il diagramma delle travature reticolari; la
	notazione dei campi fu introdotta da Robert Bow nel 1873.
\end{oss}

%\section{Esercizi proposti}
%
%\begin{esercizio}[Una coppia]
%	\label{ex:sgr_coppia}
%	Si costruiscano il poligono delle forze e un poligono funicolare del sistema
%	dell'\ESEref{ex:C02_coppia}, formato da due forze opposte di intensità $5$, parallele all'asse
%	$y$ e applicate nei punti $(1,\,0)$ e $(4,\,0)$.
%
%	\risultato Il poligono delle forze è chiuso; il primo e l'ultimo lato del poligono funicolare
%	sono paralleli, e il prodotto della loro distanza $d$ per la lunghezza $N_0$ del raggio $0P$
%	vale $15$, qualunque sia il polo: è il modulo del momento della coppia.
%\end{esercizio}
%
%\begin{esercizio}[Reazioni della trave ramificata]
%	\label{ex:sgr_ramificata}
%	Per le reazioni esterne la trave ramificata del \PARref{sec:trave_ramificata} si comporta come
%	una trave appoggiata di luce $2\ell$, soggetta a due forze verticali di intensità $p\ell$ alle
%	distanze $\ell/2$ e $3\ell/2$ dalla cerniera $A$. Si determinino le reazioni con il lato di
%	chiusura.
%
%	\risultato Il raggio parallelo al lato di chiusura incontra la linea dei carichi nel suo punto
%	medio: $V_A = Y_C = p\ell$.
%\end{esercizio}
%
%\begin{esercizio}[Ponte a stampella]
%	\label{ex:sgr_stampella}
%	Nel caso (a) dell'\ESEref{ex:ponte_stampella}, con il carico $p$ su tutta la travata, sistema
%	e carico sono simmetrici. (a) Si determini graficamente la reazione della cerniera $A$,
%	considerando il semiarco di sinistra. (b) Si scomponga la reazione secondo le due bielle che
%	concorrono in $A$.
%
%	\risultato (a) La forza trasmessa dalla cerniera $E$ è orizzontale e ha per retta d'azione
%	l'asse della travata; la risultante $\tfrac{5}{2}\,p\ell$ del carico di sinistra la incontra nel
%	punto $K$, a distanza $\tfrac{3}{4}\ell$ dalla verticale di $A$, e la reazione ha la retta $AK$,
%	di pendenza $1/3$: $V_A = \tfrac{5}{2}\,p\ell$ e $H_A = \tfrac{15}{2}\,p\ell$. (b) Il nodo $A$ è
%	un nodo canonico: il triangolo di equilibrio fornisce
%	$N_{AC} = -\tfrac{5\sqrt{5}}{12}\,p\ell$ e $N_{AD} = -\tfrac{25\sqrt{17}}{12}\,p\ell$, entrambe
%	di compressione.
%\end{esercizio}
%
%\begin{esercizio}[Diagramma cremoniano di una travatura non simmetrica]
%	\label{ex:sgr_tre_maglie}
%	Si tracci il diagramma cremoniano della travatura dell'\ESEref{ex:tre_maglie}, soggetta a una
%	forza orizzontale $2F$ in $D$ e a una forza verticale $3F$ in $E$, dopo averne determinato le
%	reazioni con la costruzione del \PARref{sec:cerniera_carrello_sgr}.
%
%	\risultato Il poligono delle forze esterne è un quadrilatero non degenere, chiuso dalle
%	reazioni $X_A = -2F$, $Y_A = \tfrac{2}{3}F$ e $Y_C = \tfrac{7}{3}F$. I nodi canonici si
%	succedono nell'ordine $A$, $C$, $D$, $E$, e il nodo $B$ fornisce la verifica. Si leggono
%	$N_{AB} = \tfrac{10}{3}F$, $N_{BC} = \tfrac{14}{3}F$, $N_{DB} = \tfrac{2\sqrt{2}}{3}F$,
%	$N_{AD} = -\tfrac{2\sqrt{5}}{3}F$, $N_{EC} = -\tfrac{7\sqrt{5}}{3}F$, $N_{DE} = -4F$ e
%	$N_{BE} = -\tfrac{2\sqrt{2}}{3}F$.
%\end{esercizio}
